%% 2026-10-15 v8：RQ1快照重算、更新實作與紀錄狀態；RQ2/RQ3仍無正式結果。
%% =====================================================================
%% 多代理人立場交換辯論機制應用於投資決策（初稿）
%% 2026-10-07 v7（最後版：補記憶規則、紀錄完整度表、一律看多檢定、RQ2/RQ3 必要對照、機制指標、放空成本限制）。
%% 2026-10-07 v6（依專題研究標準修訂：限縮結論、釐清 RQ3 定位、補記憶與稽核限制、移除未驗證占位內容，並補直接相關文獻）。
%% 2026-10-07 v5（依邏輯檢查修訂：摘要、動機銜接、模組狀態標示、第五章重排、結論補基準比較）。
%% 2026-10-07 v4（白話版：用字改為大學生程度；字距比照 v2）。v3（依 v2 仔細檢查結果修訂；改用 xelatex + xeCJK 以修正 PDF 文字層）。
%% 2026-10-06 新版研究設計修訂稿（v2：依 SSRN.tex 格式調整編號與版面）。
%% 參考文獻改由 references-meeting-20261006.bib 產生（natbib + chicago，與 SSRN.tex 相同）。
%% 編號：章 1、節 1.1、小節 1.1.1（單一阿拉伯數字層級，不再混用中文數字）。
%% 章節：摘要、引言、文獻探討、研究方法、系統設計與實作、結果與討論、結論與建議、參考文獻
%% =====================================================================

\documentclass[12pt,a4paper]{article}

%% ---------- 套件 ----------
\usepackage{amsmath}
\usepackage{fontspec}
\usepackage{xeCJK}
\setmainfont{Times New Roman}
\setsansfont{Arial}
\setCJKmainfont[Path=C:/Windows/Fonts/,AutoFakeBold=2.5]{kaiu.ttf}
\setCJKsansfont[Path=C:/Windows/Fonts/,AutoFakeBold=2.5]{kaiu.ttf}
%% 字距比照 v2（CJKutf8 + bkai 的疏朗排版）：中文字之間加一點間距
\xeCJKsetup{CJKglue=\hspace{0.1em plus 0.05em minus 0.02em}}
\usepackage[left=3cm,right=3cm,top=3cm,bottom=2cm]{geometry}
\usepackage{graphicx,array,tabularx,booktabs,float}
\usepackage{setspace,enumitem,footmisc}
\usepackage[explicit]{titlesec}
\usepackage{caption}
\usepackage{natbib}
\usepackage{tocloft}
\usepackage{tikz}
\usetikzlibrary{positioning,arrows.meta,calc}
\usepackage[bookmarks=false,hidelinks]{hyperref}

%% ---------- 行距與段落 ----------
\setstretch{1.5}
\setlength{\parindent}{2em}
\setlength{\parskip}{0pt}
\setlist{itemsep=2pt,topsep=4pt,leftmargin=2.5em}
\renewcommand{\arraystretch}{1.6}
\setlength{\tabcolsep}{8pt}
\setlength{\footnotesep}{8pt}
\emergencystretch=2em

%% ---------- 編號：章 第一章；節 第一節；小節 一、 ----------
\newcommand{\chnum}[1]{\ifcase\numexpr#1\relax 〇\or 一\or 二\or 三\or 四\or 五\or 六\or 七\or 八\or 九\or 十\or 十一\or 十二\fi}
\renewcommand{\thesection}{第\chnum{\value{section}}章}
\renewcommand{\thesubsection}{第\chnum{\value{subsection}}節}
\renewcommand{\thesubsubsection}{\chnum{\value{subsubsection}}}
\numberwithin{table}{section}
\numberwithin{figure}{section}
\renewcommand{\thetable}{\arabic{section}-\arabic{table}}
\renewcommand{\thefigure}{\arabic{section}-\arabic{figure}}
\setcounter{tocdepth}{2}
\setlength{\cftsecindent}{0pt}\setlength{\cftsecnumwidth}{4.5em}
\setlength{\cftsubsecindent}{1em}\setlength{\cftsubsecnumwidth}{4em}
\setlength{\cftsubsubsecindent}{5em}\setlength{\cftsubsubsecnumwidth}{2.5em}\renewcommand{\cftsubsubsecaftersnum}{、}
\setlength{\cfttabindent}{0pt}\renewcommand{\cfttabpresnum}{表~}\setlength{\cfttabnumwidth}{4.2em}
\setlength{\cftfigindent}{0pt}\renewcommand{\cftfigpresnum}{圖~}\setlength{\cftfignumwidth}{4.2em}
\renewcommand{\cftsecfont}{\normalfont}\renewcommand{\cftsecpagefont}{\normalfont}
\renewcommand{\cfttoctitlefont}{\hfil\fontsize{16}{24}\selectfont}\renewcommand{\cftaftertoctitle}{\hfil}
\renewcommand{\cftlottitlefont}{\hfil\fontsize{16}{24}\selectfont}\renewcommand{\cftafterlottitle}{\hfil}
\renewcommand{\cftloftitlefont}{\hfil\fontsize{16}{24}\selectfont}\renewcommand{\cftafterloftitle}{\hfil}

%% 小節（一、）的交叉參照：印出「第三章第四節第五項」
\makeatletter
\newcommand{\frlabel}[1]{\protected@write\@auxout{}{\string\global\string\@namedef{fr@#1}{\string\fr@fmt{\number\value{section}}{\number\value{subsection}}{\number\value{subsubsection}}}}}
\def\fr@fmt#1#2#3{第\chnum{#1}章第\chnum{#2}節第\chnum{#3}項}
\newcommand{\frref}[1]{\ifcsname fr@#1\endcsname\csname fr@#1\endcsname\else ??\fi}
\makeatother

%% ---------- 標題 ----------
\titleformat{\section}[block]{\centering\fontsize{18}{27}\selectfont}{\thesection}{1em}{#1}
\titleformat{\subsection}{\fontsize{16}{24}\selectfont}{\thesubsection}{1em}{#1}
\titleformat{\subsubsection}{\fontsize{14}{21}\selectfont}{\thesubsubsection、}{0em}{#1}
\titleformat{\paragraph}[runin]{\fontsize{12}{18}\selectfont}{}{0em}{#1}
\titleformat{name=\section,numberless}[block]{\centering\fontsize{18}{27}\selectfont}{}{0em}{#1}
\titleformat{name=\subsection,numberless}{\fontsize{16}{24}\selectfont}{}{0em}{#1}
\titlespacing*{\section}{0pt}{12pt}{18pt}
\titlespacing*{\subsection}{0pt}{14pt}{8pt}
\titlespacing*{\subsubsection}{0pt}{12pt}{6pt}
\titlespacing*{\paragraph}{2em}{6pt}{0.5em}

%% ---------- 圖表 ----------
\captionsetup{font={small,stretch=1.2},labelsep=space,skip=6pt}
\captionsetup[table]{position=top,justification=raggedright,singlelinecheck=false}
\setlength{\heavyrulewidth}{0.8pt}
\setlength{\lightrulewidth}{0.4pt}
\pagestyle{plain}

\newcommand{\grp}[1]{\textsf{#1}}
\newcolumntype{Y}{>{\raggedright\arraybackslash}X}
\newcolumntype{L}[1]{>{\raggedright\arraybackslash}p{#1\textwidth}}

\begin{document}

\renewcommand{\refname}{參考文獻}
\renewcommand{\contentsname}{目\quad 次}
\renewcommand{\listtablename}{表\quad 次}
\renewcommand{\listfigurename}{圖\quad 次}
\renewcommand{\tablename}{表}
\renewcommand{\figurename}{圖}

%% =====================================================================
%% 封面
%% =====================================================================
\begin{titlepage}
\centering
\setstretch{1.6}
\vspace*{2cm}
{\fontsize{18}{27}\selectfont 國立陽明交通大學管理科學系\par}
\vspace{0.3cm}
{\fontsize{16}{24}\selectfont 專題研究報告\par}
\vspace{3.5cm}
{\fontsize{20}{32}\selectfont 多代理人立場交換辯論機制\\ 應用於投資決策\par}
\vspace{0.3cm}
\vspace{0.5cm}
{\fontsize{14}{21}\selectfont Multi-Agent Role-Switching Debate\\ for Investment Decision-Making\par}
\vfill
{\fontsize{14}{24}\selectfont
\begin{tabular}{r@{\hspace{1.5em}}l}
學\hspace{2em}生： & 黃岳振 \\
指導教授： & 林斯寅\ 教授 \\
\end{tabular}\par}
\vspace{2.5cm}
{\fontsize{14}{21}\selectfont 中華民國一一五年十月\par}
\vspace{1.5cm}
\end{titlepage}

%% =====================================================================
%% 摘要
%% =====================================================================
\pagenumbering{roman}
\begin{center}
{\fontsize{16}{24}\selectfont 摘\quad 要\par}
\end{center}
\vspace{0.3cm}

投資時要同時看股價走勢、公司財報、新聞氣氛和整體經濟狀況，而這些資訊常常互相矛盾。「多代理人辯論」是讓多個 AI 角色分別扮演看多與看空，互相挑毛病，最後由一位裁決者做決定。但只是給 AI 不同的角色，並不代表它們真的會有不同的想法。本研究的核心是「自我反駁式立場交換」：讓 AI 先主張一個立場，再換邊反駁自己，最後回到原本立場整理結論。

本研究規劃三個問題：第一，整套立場交換辯論（D）比起「各守原立場」的整套固定立場辯論（C），有沒有額外的好處；第二，辯論的雙方如果使用不同的 AI 模型，這個好處會不會改變；第三，把四類資訊分給四組專長不同的辯論者處理，結果與成本會不會更好。\textbf{本報告已完成第一個問題；第二、第三個問題目前只完成研究設計，尚未執行，因此本報告不提供它們的結果。}

研究使用九檔美國大型股在 2021 至 2025 年間的 180 個季度案例，每個案例只使用分析日以前已公開的資料，並以資料快照與稽核紀錄控制主要輸入。主要評估指標是 60 個交易日後的方向準確率（猜對漲跌的比例）與年化 Sharpe 比率（每承擔一單位風險能換到多少報酬），另外也看報酬、最大虧損、雙方意見分歧與運算資源。現有版本與原始輸出的保存仍有缺漏，因此本研究將結果定位為指定樣本與流程下的回溯比較，而非可完全重現的確認性檢驗。

第一個問題的 180 案比較結果如下：D 的準確率（55.0\%）與 Sharpe（0.54）沒有比 C（58.3\%、0.59）更高，也就是沒有找到 D 比 C 好的顯著證據。此外，C、D 的準確率都低於「一律看多」的 62.2\%，四組的 Sharpe 也都低於同配置買進持有的 1.01（準確率相對「一律看多」的事後配對檢定也不顯著；Sharpe 相對買進持有則尚未檢定）。要注意的是，D 同時改變了交換立場、隔離對手、強制引用自己原文與驗證重試，所以結果不能歸因於單一的「立場交換」；而且這 180 案先前已用在試跑與規則調整，因此這是規則定好之後對既有樣本的事後回頭比較，不是對全新樣本的確認性檢驗。

\vspace{0.6cm}
\noindent 關鍵詞：大型語言模型、多代理人辯論、立場交換、模型異質性、領域專業化、投資決策

\noindent JEL 分類碼：G11、G17、C53
\clearpage
\begin{center}{\fontsize{16}{24}\selectfont Multi-Agent Role-Switching Debate\\ for Investment Decision-Making\par}\vspace{0.3cm}{\fontsize{14}{21}\selectfont Abstract\par}\end{center}
\noindent Investment decisions require the integration of technical, fundamental, news-sentiment and macroeconomic evidence, which may conflict with one another. Multi-agent debate uses opposing views and adjudication to examine arguments, yet role differences do not necessarily produce useful reasoning diversity. This study centers on self-rebuttal role-switching and plans three research questions: (RQ1) whether the full role-switching debate procedure (D) adds value over the full fixed-role debate procedure (C); (RQ2) whether model composition changes this effect; and (RQ3) whether domain-specialized decision-making improves performance and cost effectiveness. \textbf{This report completes RQ1 only; RQ2 and RQ3 are planned designs that have not been run, and no results are reported for them.}

The study uses 180 quarterly cases of nine U.S. stocks in 2021--2025, integrating four financial information domains under point-in-time data boundaries, immutable snapshots, evidence identifiers and audit trails. The primary evaluation metrics are 60-trading-day directional accuracy and annualized Sharpe ratio, supplemented by return, drawdown, disagreement and inference-resource measures.

For RQ1, D did not outperform C in directional accuracy (55.0\% vs. 58.3\%) or annualized Sharpe ratio (0.54 vs. 0.59), and no significant evidence favoring D was obtained. Point estimates for both C and D also fall below an always-buy benchmark in accuracy (62.2\%) and below a buy-and-hold reference in Sharpe ratio (1.01); a post hoc paired test of accuracy against always-buy was not significant, and the Sharpe comparison was not tested. Because D bundles role-switching, opponent isolation, forced self-quotation and validation retries, the result cannot be attributed to role-switching alone. Because the same sample had been used for pilot work and rule revisions, this is a retrospective comparison under locked rules rather than a confirmatory test.

\vspace{0.6cm}
\noindent Keywords: large language models, multi-agent debate, role switching, model heterogeneity, domain specialization, investment decisions

\noindent JEL Codes: G11, G17, C53
\clearpage
\tableofcontents
\clearpage
\listoftables
\clearpage
\listoffigures
\clearpage
\pagenumbering{arabic}

\clearpage\section{引言}

\subsection{研究背景}

\subsubsection{金融資訊整合與大型語言模型}
做投資決策時，要同時參考股價走勢（技術面）、公司財報（基本面）、新聞氣氛（情緒面）和整體經濟狀況（總體經濟面）。這些資訊更新的速度、呈現的形式，甚至指向的方向都可能不一樣。大型語言模型（例如 ChatGPT 這類 AI）可以幫忙閱讀財報、新聞和市場資料，已有不少研究把它用在金融文字分析與投資分析 \citep{lopezlira2023chatgpt,araci2019finbert,liu2023fingpt}。不過，「把資訊整理好」和「預測未來報酬」是兩件不同的事；AI 的說明寫得有條有理，不代表它真的能預測股價。

\subsubsection{多代理人辯論與立場交換}
多代理人辯論，是讓多個 AI 角色輪流提出論點、批評對方、修正自己 \citep{du2023debate,liang2023divergent}。金融領域的 AI 系統也常把工作分成分析與交易等角色 \citep{xiao2024tradingagents,zhang2024stockagent}。但辯論好不好用，要看用什麼模型、怎麼互動、怎麼裁決，並不一定比單一模型或「多問幾次再投票」更好 \citep{smit2023mad,zhang2025stop}。因此，辯論流程怎麼設計，本身就是值得研究的問題。

本研究主要比較兩種辯論方式。C 組是「固定立場辯論」：看多方一直站在看多，看空方一直站在看空。D 組是「自我反駁式立場交換」：第二輪時，每個角色只看自己第一輪寫的內容，然後換到相反立場，引用並反駁自己原本的主張；第三輪再回到原本立場，把剛剛找到的反面證據整理進結論。本研究想了解，這種「逼自己重新檢查證據」的做法，有沒有額外的價值，並進一步看它會不會受到模型組合與資訊分工方式的影響。

\subsection{研究動機}

\subsubsection{自我反駁機制有沒有額外價值}
多代理人辯論到底有沒有用，目前仍有爭議。有些研究發現辯論能提高推理的正確率 \citep{du2023debate}，但也有研究發現辯論不一定比「同一個問題問多次再投票」好，而且辯論的效果常被高估 \citep{smit2023mad,zhang2025stop}。辯論有一個常見的問題：角色會「死守立場」，只會加強自己的論點，不回應對方的證據。立場交換就是為了緩和這個問題而提出的 \citep{liu2025roleswitch,chowdhury2026proclaim}。這些研究讓我們有理由期待立場交換可能有用，但它們處理的是問答或事實查核，和預測未來股價報酬不同。所以本研究直接拿 C、D 兩種流程，在相同資料、相同預定決策呼叫次數下比較。

另外，股票回測很容易失真。AI 模型可能在訓練時就已經看過回測期間之後的股價走勢，等於「偷看答案」\citep{glasserman2023lookahead,kong2026bias}；辯論組通常也比單次判斷多用好幾倍的模型呼叫次數。因此，本專題希望在控制這些問題的前提下，檢驗立場交換在投資情境裡是否真的有用。

RQ1 的結果（見第五章）沒有找到 D 優於 C 的證據。文獻與系統設計上，有兩個條件可能限制這類機制的效果：辯論雙方由同一個模型擔任，以及四類資訊合併處理。以下兩項是\textbf{候選假說}，本研究並沒有證明它們就是 RQ1 結果的原因；RQ2 與 RQ3 的目的，是在受控條件下檢驗這兩個條件是否真的改變結果。

\subsubsection{同一個模型扮演雙方，想法可能不夠多元}
看多（Bull）與看空（Bear）兩個角色雖然被要求主張相反的立場，但如果背後是同一個 AI 模型，它們可能有相同的知識盲點、相同的推理習慣，也會犯相同的錯。角色不同，不保證想法真的不同；改用不同的模型是否能互相補足，還需要檢驗 \citep{zhang2025stop}。另一方面，不同模型的能力、訓練資料與寫作風格本來就不一樣，所以換了模型之後績效出現差異，不能直接說是「模型不同」造成的。

\subsubsection{通才組把資訊集中處理，專才組分工處理}
目前系統已有四個研究代理人，負責整理四類資訊，但真正做判斷的仍只有同一組看多與看空。技術、基本、新聞情緒與總體經濟這四類證據，更新速度和訊號性質各不相同；全部合在一起處理，可能讓某一類資訊的反面證據被稀釋，也可能讓模型一次要讀的內容太多。這是候選假說，並非已被 RQ1 的結果證實的原因。RQ3 打算讓四類資訊各自有自己的多空辯論，再由上層整合，然後評估這樣做的好處、意見衝突與額外成本。

\subsection{研究目的}

\subsubsection{立場交換機制}
檢驗「自我反駁式立場交換」相較於固定立場辯論，有沒有額外的價值。這是本研究最原始的核心問題。

\subsubsection{模型異質性}
檢驗辯論雙方使用的模型組合，會不會改變 D 相對於 C 的效果，並分辨「角色不同」與「模型不同」各自的影響。

\subsubsection{領域專業化}
檢驗「領域內辯論加上上層整合」相對於通才（Generalist）做法的績效與成本，評估把任務拆分是否划算。

\subsection{研究問題}

\subsubsection{RQ1：原始機制}
\begin{quote}在相同的案例資料與相同的預定決策呼叫次數下，自我反駁式立場交換辯論（D）的 60 日方向準確率與年化 Sharpe，是否高於固定立場辯論（C）？\end{quote}

\subsubsection{RQ2：模型異質性}
\begin{quote}當看多與看空代理人由不同的基礎模型擔任時，自我反駁式立場交換辯論（D）相對於固定立場辯論（C）的投資決策效果，是否會改變？\end{quote}

\subsubsection{RQ3：領域專業化辯論}
\begin{quote}原本是把四類金融資訊合併後，交給同一組多空代理人辯論。若改為讓技術面、基本面、情緒面與總經面代理人，各自在自己的資訊範圍內進行多空辯論，再由上層整合機制做出最終決策，方向準確率與風險調整後績效是否會更好？\end{quote}
RQ1 與 RQ2 構成主要研究線：先比較整套 C、D 流程，再檢驗模型組合是否改變兩者差異。RQ3 則是平行的系統架構延伸，檢驗資訊分工與額外運算成本的整體取捨；它不直接檢驗立場交換本身，而且比較的是「整套領域化架構」，結論不稱為單純的分工效果。

\subsection{研究範圍}

\subsubsection{樣本與資料時間邊界}
本專題是投資決策輔助的研究與系統原型，不是會自動下單的交易系統。研究範圍整理於表~\ref{tab:scope}。

\begin{table}[H]
\centering\small
\caption{研究範圍}
\label{tab:scope}
\begin{tabularx}{\textwidth}{@{}l Y@{}}
\toprule
項目 & 說明 \\
\midrule
標的與期間 & 九檔美國大型股；2021 至 2025 年每季取一個分析日，共 180 個案例 \\
資料 & 技術面、基本面、情緒面（新聞標題）與總經面；只使用分析日以前已經公開的資料 \\
決策輸出 & 買進、持有或賣出，並附上預期報酬、信心度、理由、引用的證據與風險 \\
評估 & 60 個交易日後的方向準確率與 Sharpe 比率；30 與 90 日用來檢查結果穩不穩 \\
\bottomrule
\end{tabularx}
\end{table}
\subsubsection{預測期間與決策層}

本研究維持九檔股票、20 個季末、共 180 個「股票 $\times$ 日期」案例，預測期間為 60 個交易日。C、D 兩組的候選決策只能在買進與賣出之間選一個；之後的「守門員」（Gatekeeper）可以把決策降為持有或不交易。A、B 兩組允許選持有，只作為補充對照。主要分析使用守門員之前的候選決策；經守門員處理後的實際結果，另外作為敏感性分析。以月（約 20 個交易日）為單位的比較屬於未來研究，本研究的季度樣本不變。

\subsubsection{研究狀態與證據界線}
各部分目前的進度整理於表~\ref{tab:status}。已經做完的，是在既有「季度實驗規則」下的事後回頭比較（回溯比較）。RQ2、RQ3 與月度設計都還只是設計，設計中預期的好處不能當成已有的結果。

\begin{table}[H]
\centering\small
\caption{研究狀態}
\label{tab:status}
\begin{tabularx}{\textwidth}{@{}L{.30} L{.17} Y@{}}
\toprule
項目 & 狀態 & 能說明到什麼程度 \\
\midrule
RQ1：D 對 C（qwen3:32b，季度實驗規則 v3-1002.3） & 已完成 180 案 & 這是規則定好後對既有樣本的事後比較。沒有找到 D 優於 C 的顯著證據，但這不等於兩者一樣好 \\
Gemini 3 Flash 同批比較 & 已報告（探索性） & 整套系統換成同一個別的模型，不是 RQ2；用的還是同一批案例，不算獨立重做 \\
RQ2：模型組成 & 設計階段，未執行 & 比較方式與配置規則見\frref{sec:rq2design}；尚無結果 \\
RQ3：領域分工 & 設計階段，未執行 & 流程與資源比較見\frref{sec:rq3design}；尚無結果 \\
月度／20 日設定 & 未執行 & 另外獨立設計，不併入這 180 個季度案例 \\
\bottomrule
\end{tabularx}
\end{table}

\clearpage\section{文獻探討}

\subsection{多代理人辯論}

\subsubsection{辯論與「多次取樣投票」的公平比較}
\citet{du2023debate} 讓多個模型實例進行多輪辯論後收斂出答案，使推理與事實正確性變好。\citet{liang2023divergent} 指出模型自我反省時容易死守原本的答案，因此用辯論來鼓勵多元思考。但另一方面，\citet{smit2023mad} 比較了多種辯論策略，發現多代理人辯論不一定比「自我一致性」等其他方法穩定地好，而且對設定很敏感；\citet{zhang2025stop} 則認為目前的評估常高估辯論的效果，建議重新思考評估方式，並改用不同的模型。「自我一致性」\citep{wang2023selfconsistency} 是指同一個問題問多次、再用多數決選答案。這能讓結果更穩定，也是比較辯論效果時必要的對照組；本研究的「獨立投票組」（B 組）就是依此設計。因此，本研究把重點放在流程可追蹤、比較公平（模型呼叫次數對齊），而不是先假設辯論一定比較準。
\subsection{角色交換與自我反駁}

\subsubsection{為什麼可能有效，又為什麼可能無效}
\citet{liu2025roleswitch} 提出「會考慮不確定性的角色交換辯論」：兩位辯論者先各自給出初步答案並互相質問，辯論到一半時交換角色、改為論證相反立場，然後各自回報信心與不確定性，最後由另一個裁判模型決定。在 OpenBookQA（問答題庫）上，他們的消融實驗（把某個步驟拿掉看結果變化）顯示，角色交換與回報不確定性兩個階段都有貢獻。\citet{chowdhury2026proclaim} 的 PROCLAIM 以法庭角色（原告、辯護、法官）搭配逐步檢索來查核主張，並用角色交換檢查答案是否一致；其消融結果顯示角色交換約使準確率提高 4.2 個百分點，但整體 token 用量約為標準多代理人辯論的 11 倍。\citet{chen2026mpt} 的多人格思考則讓相反的人格與中立觀點反覆對話，使偏誤被說出來並被修正。這些研究支持「從相反立場重新想一次」這個想法，但它們都不是投資決策，而且其中 PROCLAIM 的好處是靠大幅增加成本換來的，這也說明為什麼比較時必須讓模型呼叫次數相同。本研究因此採用「看多、看空、裁決者」三方架構，並讓獨立投票組與兩個辯論組的決策呼叫次數對齊（單次判斷組為基準）。\paragraph{機制的動機與可能失效的情況。}
\citet{lord1984opposite} 的人類社會判斷實驗指出，主動考慮「與自己原本信念相反的可能性」，可以減少特定的判斷偏誤。本研究據此提出假說：要求 AI 引用並反駁自己的論點，也許能挖出先前被忽略的反面證據，讓裁決者更平衡地衡量證據。不過，這個推論是從人類實驗延伸到 AI 預測股價，仍需實際檢驗；論述寫得更完整，也不一定代表預測得更準。\citet{huang2024selfcorrect} 發現，沒有外部回饋的自我修正，在其測試的推理題上可能沒效果，甚至讓表現變差。如果模型共享相同的錯誤、輸入資訊本身缺乏預測力，或市場雜訊太大，自我反駁也可能只是改變說法，績效並沒有進步。


\subsection{多代理人系統中的模型異質性}

\subsubsection{模型不同、能力不同與模型偏好}

\citet{zhang2025stop} 檢驗了多代理人辯論的評估方式與模型異質性（用不同的底層模型），讓我們有理由研究不同模型一起推理的效果。但他們的基準測試結果，不能直接套用到金融未來報酬預測，也不能直接套用到本研究的自我反駁規則。RQ2 要做的，是在同一場辯論裡同時放進 M1、M2 兩種模型；這和「整套系統都換成另一個模型」不同，後者仍然是單一模型的做法（homogeneous setting）。模型的能力高低、偏好、裁決者對某種文風的偏愛，以及不同模型犯錯的相關性，都必須分開檢查。

\citet{wang2024moa} 的 Mixture-of-Agents 把多個模型的輸出一層一層地整合，並比較「同一模型重複生成」與「多個模型各自提案」的差別。該研究也把「提出方案」與「整合方案」分成不同角色，顯示同一個模型放在流程的不同位置，適用程度可能不同。本研究據此把模型組合當成實驗要操弄的因素，並固定裁決者；不過，該研究是一般語言任務，結果不等於不同模型的金融辯論就有效，也不能排除只是模型能力有差距。

\subsection{金融領域的語言模型與投資代理人}

\subsubsection{金融文字分析與代理人流程}
\citet{lopezlira2023chatgpt} 在其樣本與設定中發現，大型語言模型對金融新聞標題的判讀，含有與後續報酬相關的訊息。FinBERT \citep{araci2019finbert} 與 FinGPT \citep{liu2023fingpt} 提供了金融文字處理的方法，FNSPID \citep{dong2024fnspid} 則提供有時間戳記的新聞資料。這些研究說明了資訊整理可能有用，但並不能直接證明任何一種辯論流程有額外價值。

TradingAgents \citep{xiao2024tradingagents} 與 StockAgent \citep{zhang2024stockagent} 把 AI 分工與交易決策放在回測或模擬環境中檢驗；它們的研究對象是整套代理人系統，和本研究「先固定案例輸入，再比較 \grp{D}、\grp{C}」的做法不同。\citet{glasserman2023lookahead} 對「偷看未來」風險的分析，以及 \citet{huang2024selfcorrect} 對「沒有外部回饋的自我修正」的謹慎結果，都提醒研究者：推理流暢、模擬績效漂亮，並不能單獨證明模型有可用的預測能力。因此，本研究把缺口限定為：在預定決策呼叫次數相同、資料時間邊界可以被查核的條件下，檢查「重新審視既有證據」的自我反駁流程，能不能改善方向準確率與風險調整後績效；而不是單純把現成的辯論方法搬到金融領域。

近期也已有研究直接檢查多代理人辯論的論證品質與金融績效關係。\citet{huang2026explainabletrading} 在歷史市場模擬中發現，較高的整體論證品質未必伴隨較高的 Sharpe 或總報酬，顯示「論證看起來更好」與「經濟結果更好」必須分開評估。該研究著重投資組合配置、論證品質與分歧維持；本研究則比較固定立場與自我反駁式立場交換的整套流程，兩者的研究對象不同，但都支持以金融績效而非文字品質作為主要判準。
\subsection{領域專業化與任務拆分}

\subsubsection{「整理資料」分工與「做決策」分工的差別}
TradingAgents 與 StockAgent 提供了金融代理人工作流程的背景 \citep{xiao2024tradingagents,zhang2024stockagent}。FinRobot 把金融任務拆給不同的專業代理人，並用分層流程處理，是金融問題拆解與工具串接的一個例子 \citep{yang2024finrobot}；但它並沒有提供本研究「四領域多空辯論 vs. 通才架構」的受控績效比較。

\citet{chen2023agentverse} 的 AgentVerse 把流程分成專家配置、協作決策、執行與評估，並討論水平與垂直兩種協作結構，說明代理人怎麼分工、怎麼溝通，會共同影響決策的過程。因此，RQ3 把專業化從「資料整理層」延伸到「領域內推理」與「上層整合」。要說明的是，把資訊切成四個領域，是本研究自己的設計選擇，並不是既有研究證明過的最佳架構。

把資訊按領域切開，可能保留每個領域的反面證據，但也可能遺失跨領域的關聯、增加摘要時壓縮資訊的損失，並帶來整合時的誤差。與其只用「模型一次讀不下太多內容」來解釋分工，本研究更關心：專業化要不要進到「決策推理」這一層？並用績效、證據使用與推論資源，來評估其中的得失。

\subsection{裁決、整合與驗證}
\subsubsection{產生論證、上層整合與證據驗證}
Mixture-of-Agents 把「提出方案」與「上層整合」分開 \citep{wang2024moa}，提供了不同角色分別負責論證與整合的架構參考。\citet{zheng2023judge} 指出，用大型語言模型當裁判時，可能出現偏好排在前面的答案、偏好較長的答案、偏好自己產生的答案等偏誤，所以裁決時輸入的順序、長度與模型標籤都需要控制。不過，該研究評估的是「回答品質」，並不能直接證明本研究的投資裁決者有預測能力。

本研究據此把「產生論證」、「整合決策」與「驗證證據」分成三個可以被查核的步驟。裁決者（Judge）或最終整合者（Final Aggregator）綜合各方論點，決定方向；守門員（Gatekeeper）則檢查證據是否真的存在、資料是否足夠，以及風險是否過高。守門員是本系統自己的驗證設計，本研究不以一般的 AI 裁判文獻，來說明它能改善投資績效。

\subsection{偷看未來的問題與金融評估}

\subsubsection{資料時間邊界與模型的訓練記憶}
\citet{glasserman2023lookahead} 指出，用 GPT 產生情緒訊號去預測報酬時，模型可能因為本來就知道這些公司的後續發展，而產生「偷看未來」的偏誤（前視偏差），並嘗試用匿名化來緩和。\citet{kong2026bias} 整理出金融語言模型研究中常見的五種偏誤（前視、存活、敘事、目標與成本偏誤），並指出多數研究沒有明確處理這些問題。因此，本研究在資料端嚴格排除分析日之後的資訊，並用前導階段做的「資料回想」與「匿名辨識」兩項探針，作為初步的風險檢查（見\frref{sec:memorydiagnostic}），回測時也扣除交易成本。不過這些探針並不是新版主樣本完整的汙染檢驗。
\subsection{文獻小結與研究缺口}

\subsubsection{從機制、模型到任務分工的連續問題}

本研究承接辯論、自我一致性、角色交換、金融文字處理與前視偏差等文獻，先從最基本的問題開始：資料固定時，自我反駁是否有額外價值？再往下檢驗，指定的模型組合與領域分工，是否會改變決策。現有文獻還不足以證明，是「模型太相似」或「資訊太多」才使 RQ1 沒看到績效進步；RQ2 與 RQ3 檢驗的正是這兩個候選假說。三個研究問題提供的是可被查核的比較，而不預設哪種架構一定比較好。此外，文獻中的問答正確率、事實查核與投資報酬預測，評估的目標並不相同，不能直接畫上等號。

